% GENERATED FILE. Edit canonical lessons, then run npm run generate:books.
% canonical-source-hash: fnv1a64:dbd4b66e1c6ade6d
% canonical-lessons: TA-C403-netuncalai, TA-C403-mempalam, TA-C403-curruvali, TA-C403-cantippu, TA-C403-virai

\chapter{Highways and Flyovers}
\label{ch:ta-highways-and-flyovers}
\hlchaptermodality{403}

\begin{chapteropening}
\textbf{By the end of this chapter:} \emph{I can talk about a highway, a flyover, a detour and a junction, and about speeding.}

Complete the last lesson of chapter 403: I can talk about a highway, a flyover, a detour and a junction, and about speeding.
\end{chapteropening}

\section[\texorpdfstring{\ta{நெடுஞ்சாலை} (neṭuñcālai)}{நெடுஞ்சாலை (neṭuñcālai)}]{\ta{நெடுஞ்சாலை} (neṭuñcālai) — a highway}

\label{lesson:TA-C403-netuncalai}



\begin{quote}
Before the new one: say the Tamil for to carry, then the Tamil for to unload.
\end{quote}

\subsection*{You'll want to know: \ta{நெடுஞ்சாலை}}
\textbf{\ta{நெடுஞ்சாலை}} — \emph{neṭuñcālai} — \textquotedblleft{}a highway\textquotedblright{}.

\begin{quote}
\ta{நெடுஞ்சாலையில்} \ta{கார்கள்} \ta{வேகமாகப்} \ta{போகின்றன}.
\end{quote}

\begin{quote}
\emph{neṭuñcālaiyil kārkaḷ vēkamākap pōkiṉṟaṉa.}
\end{quote}

\begin{quote}
Cars go fast on the highway.
\end{quote}

\ta{நெடும்}, long, + \ta{சாலை}, road: the \ta{ம்} becomes \ta{ஞ்} before \ta{ச}.

\begin{cousinweb}
\textbf{\ta{நெடுஞ்சாலை}} (\emph{neṭuñcālai}) is built from \emph{neṭu}, \textquotedblleft{}long\textquotedblright{} (DEDR 3738), and \ta{சாலை}, \textquotedblleft{}road\textquotedblright{} (DEDR 2471): the long road.
\end{cousinweb}

\begin{grammarlens}[title={What you've built}]
One more word for getting around.
\end{grammarlens}

\subsection*{Your turn}
\begin{itemize}
\raggedright
  \item \emph{Say it:} \emph{neṭuñcālai}
  \item \emph{Say it:} \emph{neṭuñcālai}, once more
  \item \emph{Say it:} \emph{neṭuñcālaiyil kārkaḷ vēkamākap pōkiṉṟaṉa.} — Cars go fast on the highway.
  \item \emph{Recall:} say the Tamil for to carry, then the Tamil for to unload, then say \emph{neṭuñcālai} again
\end{itemize}

\begin{tcolorbox}[breakable,colback=teal!4,colframe=teal!35!black,title={Before you move on}]
What does \ta{நெடுஞ்சாலை} mean? (\textquotedblleft{}A highway\textquotedblright{}.) Say it once more.
\end{tcolorbox}
\section[\texorpdfstring{\ta{மேம்பாலம்} (mēmpālam)}{மேம்பாலம் (mēmpālam)}]{\ta{மேம்பாலம்} (mēmpālam) — a flyover; an overpass}

\label{lesson:TA-C403-mempalam}



\begin{quote}
Before the new one: say the Tamil for to unload, then the Tamil for a highway.
\end{quote}

\subsection*{You'll want to know: \ta{மேம்பாலம்}}
\textbf{\ta{மேம்பாலம்}} — \emph{mēmpālam} — \textquotedblleft{}a flyover; an overpass\textquotedblright{}.

\begin{quote}
\ta{புதிய} \ta{மேம்பாலம்} \ta{நகரத்தின்} \ta{நடுவில்} \ta{இருக்கிறது}.
\end{quote}

\begin{quote}
\emph{putiya mēmpālam nakarattiṉ naṭuvil irukkiṟatu.}
\end{quote}

\begin{quote}
The new flyover is in the middle of the city.
\end{quote}

\ta{மேம்}, the joining form of \ta{மேல்}, above, + \ta{பாலம்}, a bridge.

\begin{cousinweb}
\textbf{\ta{மேம்பாலம்}} (\emph{mēmpālam}) is built from \emph{mēm}, a joining form of \emph{mēl}, \textquotedblleft{}above\textquotedblright{}, which DEDR 5086 also shows in \emph{mēm-paṭu}, \textquotedblleft{}to rise high\textquotedblright{}, and \ta{பாலம்}, \textquotedblleft{}a bridge\textquotedblright{} (DEDR 4099): a bridge over the road.
\end{cousinweb}

\begin{grammarlens}[title={What you've built}]
One more word for getting around.
\end{grammarlens}

\subsection*{Your turn}
\begin{itemize}
\raggedright
  \item \emph{Say it:} \emph{mēmpālam}
  \item \emph{Say it:} \emph{mēmpālam}, once more
  \item \emph{Say it:} \emph{putiya mēmpālam nakarattiṉ naṭuvil irukkiṟatu.} — The new flyover is in the middle of the city.
  \item \emph{Recall:} say the Tamil for to unload, then the Tamil for a highway, then say \emph{mēmpālam} again
\end{itemize}

\begin{tcolorbox}[breakable,colback=teal!4,colframe=teal!35!black,title={Before you move on}]
What does \ta{மேம்பாலம்} mean? (\textquotedblleft{}A flyover; an overpass\textquotedblright{}.) Say it once more.
\end{tcolorbox}
\section[\texorpdfstring{\ta{சுற்றுவழி} (cuṟṟuvaḻi)}{சுற்றுவழி (cuṟṟuvaḻi)}]{\ta{சுற்றுவழி} (cuṟṟuvaḻi) — a detour; a roundabout way}

\label{lesson:TA-C403-curruvali}



\begin{quote}
Before the new one: say the Tamil for a highway, then the Tamil for a flyover.
\end{quote}

\subsection*{You'll want to know: \ta{சுற்றுவழி}}
\textbf{\ta{சுற்றுவழி}} — \emph{cuṟṟuvaḻi} — \textquotedblleft{}a detour; a roundabout way\textquotedblright{}.

\begin{quote}
\ta{பாலம்} \ta{மூடியிருந்ததால்} \ta{சுற்றுவழியில்} \ta{போனோம்}.
\end{quote}

\begin{quote}
\emph{pālam mūṭiyiruntatāl cuṟṟuvaḻiyil pōṉōm.}
\end{quote}

\begin{quote}
Because the bridge was closed, we went by a detour.
\end{quote}

\ta{சுற்று}, to go around, + \ta{வழி}, way: the long way round.

\begin{cousinweb}
\textbf{\ta{சுற்றுவழி}} (\emph{cuṟṟuvaḻi}) is built from \ta{சுற்று}, \textquotedblleft{}to go around\textquotedblright{} (DEDR 2715), and \ta{வழி}, \textquotedblleft{}a way\textquotedblright{} (DEDR 5297).
\end{cousinweb}

\begin{grammarlens}[title={What you've built}]
One more word for getting around.
\end{grammarlens}

\subsection*{Your turn}
\begin{itemize}
\raggedright
  \item \emph{Say it:} \emph{cuṟṟuvaḻi}
  \item \emph{Say it:} \emph{cuṟṟuvaḻi}, once more
  \item \emph{Say it:} \emph{pālam mūṭiyiruntatāl cuṟṟuvaḻiyil pōṉōm.} — Because the bridge was closed, we went by a detour.
  \item \emph{Recall:} say the Tamil for a highway, then the Tamil for a flyover, then say \emph{cuṟṟuvaḻi} again
\end{itemize}

\begin{tcolorbox}[breakable,colback=teal!4,colframe=teal!35!black,title={Before you move on}]
What does \ta{சுற்றுவழி} mean? (\textquotedblleft{}A detour; a roundabout way\textquotedblright{}.) Say it once more.
\end{tcolorbox}
\section[\texorpdfstring{\ta{சந்திப்பு} (cantippu)}{சந்திப்பு (cantippu)}]{\ta{சந்திப்பு} (cantippu) — a junction; a meeting}

\label{lesson:TA-C403-cantippu}



\begin{quote}
Before the new one: say the Tamil for a flyover, then the Tamil for a detour.
\end{quote}

\subsection*{You'll want to know: \ta{சந்திப்பு}}
\textbf{\ta{சந்திப்பு}} — \emph{cantippu} — \textquotedblleft{}a junction; a meeting\textquotedblright{}.

\begin{quote}
\ta{அந்தச்} \ta{சந்திப்பில்} \ta{பேருந்து} \ta{இடது} \ta{பக்கம்} \ta{திரும்பியது}.
\end{quote}

\begin{quote}
\emph{antac cantippil pēruntu iṭatu pakkam tirumpiyatu.}
\end{quote}

\begin{quote}
At that junction the bus turned left.
\end{quote}

\ta{ரயில்} \ta{சந்திப்பு} is a railway junction; the same word names a meeting of people.

\begin{cousinweb}
\textbf{\ta{சந்திப்பு}} (\emph{cantippu}) is built from \ta{சந்தி}, \textquotedblleft{}to meet\textquotedblright{}, borrowed from Sanskrit \emph{saṁdhi}, \textquotedblleft{}a joining, a joint\textquotedblright{} (CDIAL 12913), with the noun ending \emph{-ppu}.
\end{cousinweb}

\begin{grammarlens}[title={What you've built}]
One more word for getting around.
\end{grammarlens}

\subsection*{Your turn}
\begin{itemize}
\raggedright
  \item \emph{Say it:} \emph{cantippu}
  \item \emph{Say it:} \emph{cantippu}, once more
  \item \emph{Say it:} \emph{antac cantippil pēruntu iṭatu pakkam tirumpiyatu.} — At that junction the bus turned left.
  \item \emph{Recall:} say the Tamil for a flyover, then the Tamil for a detour, then say \emph{cantippu} again
\end{itemize}

\begin{tcolorbox}[breakable,colback=teal!4,colframe=teal!35!black,title={Before you move on}]
What does \ta{சந்திப்பு} mean? (\textquotedblleft{}A junction; a meeting\textquotedblright{}.) Say it once more.
\end{tcolorbox}
\section[\texorpdfstring{\ta{விரை} (virai)}{விரை (virai)}]{\ta{விரை} (virai) — to speed; to hurry}

\label{lesson:TA-C403-virai}



\begin{quote}
Before the new one: say the Tamil for a detour, then the Tamil for a junction.
\end{quote}

\subsection*{You'll want to know: \ta{விரை}}
\textbf{\ta{விரை}} — \emph{virai} — \textquotedblleft{}to speed; to hurry\textquotedblright{}.

\begin{quote}
\ta{செய்தி} \ta{கேட்டதும்} \ta{காவலர்கள்} \ta{அந்த} \ta{இடத்துக்கு} \ta{விரைந்தார்கள்}.
\end{quote}

\begin{quote}
\emph{ceyti kēṭṭatum kāvalarkaḷ anta iṭattukku viraintārkaḷ.}
\end{quote}

\begin{quote}
As soon as they heard the news, the police hurried to that place.
\end{quote}

Past (I …): \ta{விரைந்தேன்}, \emph{viraintēṉ}.

Mostly written and in the news; in speech, \ta{வேகமாகப்} \ta{போ}.

\begin{cousinweb}
\textbf{\ta{விரை}} (\emph{virai}) is native Dravidian (DEDR 5417), \textquotedblleft{}to be swift, to hurry\textquotedblright{}.
\end{cousinweb}

\begin{grammarlens}[title={What you've built}]
One more word for getting around.
\end{grammarlens}

\subsection*{Your turn}
\begin{itemize}
\raggedright
  \item \emph{Say it:} \emph{virai}
  \item \emph{Say it:} \emph{virai}, once more
  \item \emph{Say it:} \emph{ceyti kēṭṭatum kāvalarkaḷ anta iṭattukku viraintārkaḷ.} — As soon as they heard the news, the police hurried to that place.
  \item \emph{Recall:} say the Tamil for a detour, then the Tamil for a junction, then say \emph{virai} again
\end{itemize}

\begin{tcolorbox}[breakable,colback=teal!4,colframe=teal!35!black,title={Before you move on}]
What does \ta{விரை} mean? (\textquotedblleft{}To speed; to hurry\textquotedblright{}.) Say it once more.
\end{tcolorbox}
